%=====================================================================
%  Chapter 4 — Embedded Firmware: the ESP32 micro-ROS Node
%=====================================================================
\chapter{Embedded Firmware: the On-Board Controller}
\label{chap:firmware}

\section{Introduction}
The firmware is the lowest software layer of the robot: it runs directly
on the microcontroller introduced in Chapter~\ref{chap:hardware} and
turns the electronics into an active component of the robot's software
system. Its responsibilities are concrete: receive the eighteen target
joint angles, drive the actuated servos smoothly through the PCA9685
board, read the MPU6050 inertial sensor, and make all of this available to the
rest of the system. This chapter first justifies the choice of the ESP32,
then explains how the firmware is organised, how it handles the actuators
and sensors over the I\textsuperscript{2}C bus, and finally how it joins
the ROS\,2 world through \emph{micro-ROS} --- the bridge that leads
naturally into the next chapter.

\section{Choice of the ESP32}
\label{sec:choice-esp}
The on-board controller has a demanding but bounded role: it must update
eighteen servos and sample an inertial sensor at a steady rate, while
exchanging data with the host computer. Selecting a microcontroller for
this task involves balancing computing power, connectivity, and cost.

The \textbf{ESP32} was chosen because it meets all of these constraints
at once. It is a low-cost, dual-core processor running at 240~MHz with
enough memory to host a lightweight ROS\,2 client, which is far beyond
the reach of an 8-bit board such as a basic Arduino. Crucially, it
integrates \textbf{WiFi}, allowing the robot to communicate with the host
wirelessly and to operate untethered. It natively supports the
\textbf{I\textsuperscript{2}C} bus required by both the PCA9685 PWM
drivers and the MPU6050 sensor, and it runs a real-time operating system
(FreeRTOS) whose two cores can separate time-critical communication from
supervision. Finally, it is well supported by the embedded ROS\,2
tooling used in this project. Compared with a single-board computer such
as a Raspberry Pi, the ESP32 is cheaper, lighter, and consumes far less
power --- and the modest computation performed on board does not require
a full operating system, making the single-board computer unnecessary
for this layer.

\section{Firmware Architecture}
\label{sec:fw-arch}
The firmware is written in C++ on the Arduino/PlatformIO framework and is
deliberately split into focused modules rather than a single monolithic
program: one module manages the servos, one is a self-contained driver
for the IMU, and one handles the communication with the host. A separate
file holds the joint map (Table~\ref{tab:joints}) and all hardware
constants, so that no hardware-specific value is hidden in the logic.

The architecture takes advantage of the ESP32's two cores. All
communication --- the wireless link, the connection lifecycle, and the
dispatching of incoming and outgoing messages --- runs in a dedicated
real-time task pinned to \textbf{core~1}. \textbf{Core~0} is left to the
main program, which only supervises the system and is reserved for future
watchdog and safety logic. Separating communication from supervision in
this way ensures that activity on the network cannot stall the safety
logic, and vice versa.

Functionally, the firmware implements two continuous data paths, shown in
Figure~\ref{fig:fwflow}. On the \emph{command path}, target joint angles
arriving from the host are stored, and a periodic 100~Hz routine drives
the servos toward those targets through the PWM boards. On the
\emph{sensing path}, the inertial sensor is sampled and a periodic 50~Hz
routine sends the measurements back to the host. Both periodic routines
are driven by the microcontroller's own clock, so the joint-control rate
remains regular even when network traffic is not. On start-up, the same
servo mechanism first brings all joints smoothly to their resting
position before any command is accepted.

\begin{figure}[ht]
  \centering
  \begin{tikzpicture}[
      font=\small, >=Latex,
      mod/.style={draw, rounded corners=2pt, align=center, thick,
                  minimum height=9mm, fill=orange!10},
      hw/.style={draw, rounded corners=2pt, align=center, thick,
                 minimum height=9mm, fill=gray!12},
      ext/.style={draw, rounded corners=2pt, align=center, thick,
                  minimum height=9mm, fill=blue!8},
    ]
    \node[ext] (host) {ROS\,2\\host};
    \node[mod, right=20mm of host, yshift=11mm] (sub) {Command\\handler};
    \node[mod, right=12mm of sub] (servo) {Servo loop\\100\,Hz};
    \node[mod, below=12mm of sub] (imudrv) {IMU driver\\(I\textsuperscript{2}C)};
    \node[mod, right=12mm of imudrv] (imutim) {IMU loop\\50\,Hz};
    \node[hw, right=14mm of servo] (pca) {PCA9685\\$\rightarrow$ 16 servos};
    \node[hw, right=14mm of imutim] (imu) {MPU6050};

    \begin{scope}[on background layer]
      \node[draw, dashed, thick, rounded corners, fill=green!4,
            fit=(sub)(servo)(imudrv)(imutim), inner sep=10pt] (esp) {};
      \node[anchor=south west, font=\scriptsize\itshape]
            at (esp.north west) {ESP32 (FreeRTOS, core 1)};
    \end{scope}

    \draw[->] (host.north) |- (sub.west) node[midway,above,font=\scriptsize]{};
    \draw[->] (sub) -- (servo);
    \draw[->] (servo) -- node[above,font=\scriptsize]{I\textsuperscript{2}C} (pca);
    \draw[->] (imu) -- node[above,font=\scriptsize]{I\textsuperscript{2}C} (imutim);
    \draw[->] (imutim) -- (imudrv);
    \draw[->] (imudrv.west) -| (host.south);
    \node[font=\scriptsize] at ($(host)!0.5!(sub)+(0,4mm)$) {/joint\_commands};
    \node[font=\scriptsize] at ($(host)!0.5!(imudrv)+(2mm,-4mm)$) {/imu};
  \end{tikzpicture}
  \caption{Firmware data paths: the command path (top) drives the servos
           at 100~Hz, the sensing path (bottom) publishes the IMU at
           50~Hz. Both periodic loops run on the microcontroller's clock.}
  \label{fig:fwflow}
\end{figure}

The on-board electronics that implement this architecture are shown in
Figure~\ref{fig:onboard}: the ESP32 controller with its servo-driver
board and wiring harness, and the DC-DC regulators that power the
actuators.

\begin{figure}[ht]
  \centering
  \begin{subfigure}{0.52\textwidth}
    \centering
    \includegraphics[height=5.5cm]{bor1.jpeg}
    \caption{ESP32 controller and PWM servo driver}
    \label{fig:onboard-esp}
  \end{subfigure}
  \hfill
  \begin{subfigure}{0.40\textwidth}
    \centering
    \includegraphics[height=5.5cm]{bor2.jpeg}
    \caption{DC-DC power regulators}
    \label{fig:onboard-pwr}
  \end{subfigure}
  \caption{On-board electronics of Abbes: (a) the ESP32 with the PCA9685
           servo driver and the servo wiring harness, and (b) the DC-DC
           regulators supplying the actuators.}
  \label{fig:onboard}
\end{figure}

\section{Actuators and Sensors over I\texorpdfstring{\textsuperscript{2}}{2}C}
\label{sec:fw-io}
The PWM driver boards and the inertial sensor all share a single
\textbf{I\textsuperscript{2}C} bus, a two-wire serial bus on which each
device has its own address. This is what allows the peripherals --- the
PCA9685 PWM driver (at address 0x40) and the MPU6050 --- to be controlled
through only two microcontroller pins.

\paragraph{Actuators.}
A joint angle cannot be sent to a servo directly; it must be expressed as
the width of a PWM pulse. The firmware converts each angle in degrees to
the corresponding pulse width for the MG996R servos and sends it to the
appropriate PCA9685 channel. Of the eighteen joints in the command
vector, sixteen are wired to the board's channels (the ten leg joints and
the six arm joints); the two head joints are still accepted in the
command but have no servo in the current build and are simply skipped.
Two safeguards are applied in the process.
First, every command is \emph{clamped} to the individual angular limits
of the joint (Table~\ref{tab:joints}) before it reaches a servo, so an
out-of-range request can never force a joint past its mechanical limit.
Second, and most importantly, the servos are never commanded to jump
straight to a new target: the target is only stored, and the 100~Hz loop
moves each joint toward it by a small, fixed maximum step on every tick
--- a \emph{velocity-limited interpolation}. This ramps every motion
smoothly and protects the plastic gears of the hobby servos from the
abrupt slams that a sudden set-point would otherwise cause.

\paragraph{Sensor.}
The MPU6050 combines a three-axis accelerometer and a three-axis
gyroscope, and is read over the same I\textsuperscript{2}C bus through a
self-contained driver rather than a third-party library, keeping the
behaviour fully under control. At start-up the driver wakes the sensor,
enables a low-pass filter to reject vibration noise, and sets the
measurement ranges and sampling rate. Each reading then retrieves all of
the acceleration and angular-velocity data in a single bus transaction,
which guarantees that every axis belongs to the same sample, and converts
the raw values into physical units. Importantly, the firmware performs
\emph{no orientation estimation on the device}: it transmits acceleration
and angular velocity only. Fusing those two signals into a usable tilt
estimate is left to a node on the host side
(Chapter~\ref{chap:software}), which keeps the microcontroller's workload
light and places the filtering logic where it is easiest to tune.

\section{Joining ROS\,2 with micro-ROS}
\label{sec:fw-microros}
The two data paths described above must ultimately connect to the rest of
the system. The entire software stack of Abbes is built on ROS\,2 --- the
simulation, the controllers, and the artificial-intelligence layer of the
following chapters all communicate through it. ROS\,2, however, is
designed to run on full computers: its communication relies on the DDS
middleware, whose memory and processing requirements far exceed what a
microcontroller can provide. Keeping the ESP32 outside ROS\,2 and talking
to it over a custom serial protocol would mean hand-written message
framing on both sides, no introspection, and brittle, hard-to-extend
code.

\textbf{micro-ROS} resolves this. It brings the ROS\,2 client concepts
--- nodes, topics, publishers and subscribers --- to microcontrollers
running a real-time operating system. In place of the full DDS middleware
it uses a lightweight client that communicates with a companion process
on the host, the \textbf{micro-ROS agent}; the agent translates the
compact messages from the microcontroller into the full ROS\,2 traffic,
and back (Figure~\ref{fig:microros}). From the point of view of the rest
of the system, the ESP32 then appears as an ordinary ROS\,2 node: its
data can be inspected with the same tools and uses the same standard
message types as any other node.

\begin{figure}[ht]
  \centering
  \begin{tikzpicture}[
      font=\small, >=Latex, node distance=20mm,
      b/.style={draw, rounded corners=2pt, align=center, thick,
                minimum height=12mm, minimum width=30mm},
    ]
    \node[b, fill=orange!12] (esp) {ESP32\\\scriptsize micro-ROS client};
    \node[b, fill=green!10, right=24mm of esp] (agent) {micro-ROS agent\\\scriptsize (host)};
    \node[b, fill=blue!8, right=24mm of agent] (ros) {ROS\,2 graph\\\scriptsize controllers, sim, AI};
    \draw[<->] (esp) -- node[above,font=\scriptsize]{WiFi/UDP} (agent);
    \draw[<->] (agent) -- node[above,font=\scriptsize]{DDS} (ros);
  \end{tikzpicture}
  \caption{micro-ROS bridges the microcontroller to the ROS\,2 graph: a
           lightweight client on the ESP32 talks to an agent on the host,
           which connects it to the full DDS network.}
  \label{fig:microros}
\end{figure}

This is exactly why micro-ROS is convenient here: the on-robot controller
speaks the same language as everything else, so the inertial data and the
joint commands are exposed as standard ROS\,2 topics, with no custom
protocol or translation glue to maintain. Concretely, the firmware
advertises a node that \emph{subscribes} to the joint-command topic and
\emph{publishes} the IMU topic. Both use \textbf{best-effort} delivery:
for a continuous stream of set-points and inertial samples, the newest
value always supersedes the previous one, so re-transmitting a lost
message would only add latency.

Because the wireless link can drop, the firmware tolerates an agent that
comes and goes through the explicit state machine of
Figure~\ref{fig:fsm}. It waits for the network, then for the agent; once
the agent answers it creates the ROS\,2 entities and enters the connected
state, where it periodically re-checks the link. If the agent is lost,
the node cleanly destroys its entities and returns to waiting --- so the
robot reconnects automatically, without a reboot. An on-board LED mirrors
this state, lit only while connected.

\begin{figure}[ht]
  \centering
  \begin{tikzpicture}[
      font=\small, >=Latex, node distance=18mm and 24mm,
      st/.style={draw, rounded corners=3pt, align=center, thick,
                 minimum height=9mm, minimum width=26mm, fill=blue!7}
    ]
    \node[st] (wifi) {\textsc{Waiting\_WiFi}};
    \node[st, right=of wifi] (agent) {\textsc{Waiting\_Agent}};
    \node[st, right=of agent] (conn) {\textsc{Connected}};
    \node[st, below=of agent] (disc) {\textsc{Disconnected}};

    \draw[->] (wifi) -- node[above,font=\scriptsize,align=center]{WiFi up} (agent);
    \draw[->] (agent) -- node[above,font=\scriptsize,align=center]{agent replies} (conn);
    \draw[->] (conn) -- node[right,font=\scriptsize,align=center]{link lost} (disc);
    \draw[->] (disc) -- node[below,font=\scriptsize,align=center]{entities destroyed} (agent);
  \end{tikzpicture}
  \caption{Agent-connection state machine, providing automatic
           reconnection without rebooting the robot.}
  \label{fig:fsm}
\end{figure}

\section{Conclusion}
This chapter described the embedded firmware of Abbes. The ESP32 was
chosen for its dual-core processing, integrated WiFi, native
I\textsuperscript{2}C support, and low cost. Around it, the firmware is
organised into focused modules running two periodic data paths --- a
100~Hz servo loop and a 50~Hz sensing loop --- with per-joint limiting
and velocity-limited interpolation to protect the hardware. The actuators
and the inertial sensor share a single I\textsuperscript{2}C bus, and the
node joins the rest of the system through micro-ROS, which makes the
microcontroller a first-class participant in the ROS\,2 graph. That graph
--- the host-side software architecture that produces the joint commands
and consumes the inertial data --- is the subject of the next chapter.
